\documentclass[11pt]{article}
\usepackage[margin=1.1in]{geometry}
\usepackage{amsmath,amssymb,amsthm}
\usepackage{booktabs}
\usepackage[hidelinks]{hyperref}
\usepackage{orcidlink}
\usepackage{xcolor}

\newtheorem{theorem}{Theorem}
\newtheorem{proposition}{Proposition}
\newtheorem{lemma}{Lemma}
\newtheorem{corollary}{Corollary}
\newtheorem{definition}{Definition}
\newtheorem*{namedconj}{Conjecture (Goemans)}
\newcommand{\doi}[1]{\href{https://doi.org/#1}{doi:#1}}
\theoremstyle{remark}
\newtheorem{remark}{Remark}

\newcommand{\dmax}{d_{\max}}
\newcommand{\flow}{f}

% ---------------------------------------------------------------------------
% Standard CAOS Research manuscript header block. Every manuscript states, on
% page 1 and without the reader having to infer it: what KIND of document it is,
% which programme it belongs to, its version, its licence, and its DOI.
% ---------------------------------------------------------------------------
\newcommand{\doctype}{Preprint}
\newcommand{\docsubtype}{independent verification and replication record}
\newcommand{\docversion}{v0.01}
\newcommand{\docdate}{25 July 2026}
\newcommand{\doclicense}{CC BY 4.0}
\newcommand{\docdoi}{10.5281/zenodo.21554259}
\newcommand{\docconceptdoi}{10.5281/zenodo.21554258}
\newcommand{\docrepo}{github.com/fsantibanezleal/CAOS\_RESEARCH}

\newcommand{\headerblock}{%
\begin{center}
\fbox{\parbox{0.94\textwidth}{\small
\textbf{\doctype} (\docsubtype) \quad$\cdot$\quad \docversion, \docdate \quad$\cdot$\quad
Licence \doclicense\\[2pt]
CAOS Research open-problems programme \quad$\cdot$\quad problem
\texttt{unsplittable-flow-cost}, experiments 001--003\\[2pt]
Version DOI \href{https://doi.org/\docdoi}{\texttt{\docdoi}} \quad$\cdot$\quad
concept DOI (always latest) \href{https://doi.org/\docconceptdoi}{\texttt{\docconceptdoi}}\\[2pt]
Not peer reviewed. Source, declared hypotheses, artifacts and verdicts:
\texttt{\docrepo}
}}
\end{center}
\vspace{4pt}}

\title{An independent exact verification of the 2026 counterexample to
Goemans' unsplittable-flow cost conjecture, with the violation constant it
forces\\
{\large Machine record v0.01: the CAOS Research unsplittable-flow-cost program,
experiments 001--003}}
\author{Felipe Santiba\~nez-Leal\orcidlink{0000-0002-0150-3246}%
\thanks{CAOS open-research programme, Santiago, Chile. With machine-executed
computations; all code, declared hypotheses, artifacts and verdicts are
versioned at \url{https://github.com/fsantibanezleal/CAOS_RESEARCH} under
\texttt{problems/optimization-geometry/unsplittable-flow-cost/}.}}
\date{July 25, 2026}

\begin{document}
\maketitle
\headerblock

\begin{abstract}
In July 2026 a counterexample to Goemans' cost version of the
Dinitz--Garg--Goemans theorem was announced publicly, outside peer review, and
circulated widely. We report an independent verification of that finite object
by exact rational enumeration, performed without executing or importing the
proposer's verifier, together with quantitative and structural content that the
announcement did not state. Our results: (i) the instance is confirmed, so
Goemans' conjecture is false, and, the instance being acyclic, the
Morell--Skutella cost conjecture and the convex-combination form fall with it;
(ii) the instance forces a violation of exactly $\tfrac{16}{15}\dmax$, that is,
the conjecture fails by ONE unit at $\dmax = 15$, so the open question whether
cost-preserving rounding is possible with $O(\dmax)$ violation is untouched by
it; (iii) every proved result in the area survives its direct test on the
instance, including the Dinitz--Garg--Goemans theorem, the Morell--Skutella
cost-free conjecture, the series-parallel theorem of Majthoub Almoghrabi,
Skutella and Warode, and the planar theorem of Traub, Vargas Koch and
Zenklusen; (iv) the instance contains a $K_4$ subdivision, hence sits exactly
one structure outside the series-parallel class where the conjecture is proved,
and, being planar, pins the planar constant strictly between 1 and 2; (v) an
exact separation linear program shows that the published cost vector is an
optimal separator for that graph and fractional flow, so its cost gap cannot be
improved by any reweighting; and (vi) a first minimality result, that no
single-terminal instance is a counterexample for any nonnegative cost vector,
together with a necessary condition at two terminals. The verification is a
replication; items (ii) and (iv)--(vi) are, to our knowledge, recorded here for
the first time. A prediction of ours about two terminals was refuted by our own
machine and is retracted in Section~\ref{sec:minimality}.
\end{abstract}

\section{The statement}

Let $G = (V, A)$ be a digraph with a source $s$ and terminals $T$ carrying
demands $d \in \mathbb{Q}_{\geq 0}^T$, and let $x \in \mathbb{Q}_{\geq 0}^A$ be
a \emph{fractional flow} routing $d_t$ to each $t \in T$. An \emph{unsplittable
flow} $\mathcal{P} = \{P^t\}_{t \in T}$ assigns one $s$--$t$ path per terminal,
with arc load
\[
\flow^{\mathcal{P}}(a) \;=\; \sum_{t \in T \colon a \in P^t} d_t ,
\qquad \dmax \;=\; \max_{t \in T} d_t .
\]
Throughout this literature the capacity of an arc is taken to be $x(a)$ itself,
so violation is measured against the flow being rounded \cite{TVZ24}.

\begin{theorem}[Dinitz, Garg, Goemans \cite{DGG99}]
Every such instance admits, in polynomial time, an unsplittable flow with
$\flow^{\mathcal{P}}(a) \le x(a) + \dmax$ for every arc $a$.
\end{theorem}

Call such a routing \emph{congestion-good}, and call it \emph{cost-good} for a
cost vector $c \in \mathbb{Q}^A_{\geq 0}$ if
$c^{\mathsf T}\flow^{\mathcal{P}} \le c^{\mathsf T} x$.

\begin{namedconj}
For every instance and every $c \ge 0$ there is an unsplittable flow that is
simultaneously congestion-good and cost-good.
\end{namedconj}

Each condition alone is easy or known: congestion-goodness alone is the theorem
above, and cost-goodness alone is achieved by routing every terminal on a
cheapest path. The content of the conjecture is the simultaneity. Its status
before 2026: proved when all demands are multiples of one another
\cite{Sku02}; proved for series-parallel digraphs, in the stronger
convex-combination form and with strict deviation below $\dmax$ \cite{MSW25};
and known for planar graphs only at the doubled violation $2\dmax$
\cite{TVZ24}. As of October 2025 the primary literature recorded that whether
the conjecture holds even with $O(\dmax)$ violation ``remains wide open'' and
that proving such a weaker form ``would already be considered a breakthrough''
\cite{STVZ25}.

\section{The instance}

Vertices $\{s, u, v, w, t_1, t_2, t_3\}$, demands $d = (15, 10, 15)$ so
$\dmax = 15$, and nine arcs:

\begin{center}
\begin{tabular}{@{}lrrr@{}}
\toprule
arc $a$ & $x_a$ & $c_a$ & $x_a + \dmax$\\
\midrule
$s \to t_1$ & 10 & 2 & 25\\
$s \to t_2$ & 6 & 3 & 21\\
$s \to u$ & 24 & 0 & 39\\
$u \to t_3$ & 10 & 2 & 25\\
$u \to v$ & 14 & 0 & 29\\
$v \to t_1$ & 5 & 0 & 20\\
$v \to w$ & 9 & 0 & 24\\
$w \to t_2$ & 4 & 0 & 19\\
$w \to t_3$ & 5 & 0 & 20\\
\bottomrule
\end{tabular}
\end{center}

The flow is feasible (conservation verified at all seven vertices) with
$c^{\mathsf T} x = 2\cdot 10 + 3 \cdot 6 + 2 \cdot 10 = 58$.

\section{Verification}

\paragraph{Method and independence.} We wrote an exact checker from the
conjecture statement as printed in \cite{STVZ25}. It models arcs by index, so
parallel arcs are representable; it enumerates simple $s$--$t$ paths by its own
depth-first search, never from a supplied list; it decides both conditions with
inclusive inequalities in \texttt{fractions.Fraction} arithmetic, with floats
refused by the instance constructor. It was calibrated first on instances whose
answers are fixed by inspection, including a parallel-arc instance, an instance
with a directed cycle, and an instance engineered to load an arc to exactly
$x_a + \dmax$. The proposer's verifier was archived and hashed but never
imported or executed, and the instance was re-entered by hand rather than
parsed from their file, so that agreement is evidence rather than a shared-code
artifact. The historical failure mode of candidate counterexamples in this area
is incomplete path enumeration, so the path count is asserted, and the
Dinitz--Garg--Goemans theorem is used as a per-instance oracle: an enumeration
reporting no congestion-good routing would indicate an error in our data or
code, never in the theorem.

\paragraph{Result.} Each terminal has exactly two simple paths, an expensive
direct choice $E_i$ and a free detour $Z_i$:
\[
E_1 = s t_1,\quad Z_1 = s\,u\,v\,t_1; \qquad
E_2 = s t_2,\quad Z_2 = s\,u\,v\,w\,t_2; \qquad
E_3 = s\,u\,t_3,\quad Z_3 = s\,u\,v\,w\,t_3 ,
\]
each $E_i$ costing exactly 30 when carrying its demand and each $Z_i$ costing
zero. All eight routings:

\begin{center}
\begin{tabular}{@{}ccccclll@{}}
\toprule
$t_1$ & $t_2$ & $t_3$ & cost & $\alpha$ & status & violated arcs\\
\midrule
$E_1$ & $E_2$ & $E_3$ & 90 & $1/3$ & congestion-good & --\\
$E_1$ & $E_2$ & $Z_3$ & 60 & $2/3$ & congestion-good & --\\
$E_1$ & $Z_2$ & $E_3$ & 60 & $2/5$ & congestion-good & --\\
$E_1$ & $Z_2$ & $Z_3$ & 30 & $16/15$ & cost-good only & $v \to w$ by 1\\
$Z_1$ & $E_2$ & $E_3$ & 60 & $2/3$ & congestion-good & --\\
$Z_1$ & $E_2$ & $Z_3$ & 30 & $16/15$ & cost-good only & $u \to v$ by 1\\
$Z_1$ & $Z_2$ & $E_3$ & 30 & $16/15$ & cost-good only & $s \to u$ by 1\\
$Z_1$ & $Z_2$ & $Z_3$ & 0 & $26/15$ & cost-good only & $s\to u$ by 1, $u \to v$ by 11, $v \to w$ by 1\\
\bottomrule
\end{tabular}
\end{center}

Here $\alpha$ denotes the least budget, in units of $\dmax$, at which the
routing becomes congestion-good. The congestion-good and cost-good sets are
disjoint.

\begin{theorem}[machine-verified]
For the instance above, no unsplittable flow satisfies both
$\flow^{\mathcal{P}}(a) \le x(a) + \dmax$ for all $a$ and
$c^{\mathsf T}\flow^{\mathcal{P}} \le c^{\mathsf T} x$. Goemans' conjecture is
false.
\end{theorem}

\paragraph{An independent second route.} The same conclusion follows without
costing the eight routings. Call two free choices \emph{in conflict} when every
completion of the remaining terminal violates some arc bound. Then $Z_1 Z_3$
conflict ($u \to v$ carries $30 > 29$), $Z_2 Z_3$ conflict ($v \to w$ carries
$25 > 24$), and $Z_1 Z_2$ conflict ($s \to u$ carries $40 > 39$, because $t_3$
enters through $s \to u$ on either of its paths). The conflict graph is a
triangle, its independence number is 1, so every congestion-good routing pays
for at least two expensive choices and costs at least $30 + 30 = 60 > 58$. Our
implementation computes the conflict graph and this bound in code independent of
the enumeration, and the two agree.

\paragraph{The mechanism is a stable-set gap.} Writing $\rho_i$ for the fraction
of terminal $i$'s demand that the fractional flow sends on its free choice,
\[
\rho \;=\; \left(\tfrac{5}{15}, \tfrac{4}{10}, \tfrac{5}{15}\right)
\;=\; \left(\tfrac13, \tfrac25, \tfrac13\right), \qquad
\textstyle\sum_i \rho_i \;=\; \tfrac{16}{15} \;>\; 1 .
\]
A congestion-good routing selects an independent set of the conflict graph and
so buys at most one unit of free routing, while the fractional flow buys
$16/15$. The instance is the linear-programming integrality gap of the
stable-set polytope on a triangle, transported into flow language, with the arc
costs acting as a nonnegative separator of the violated triangle inequality.

\section{How much is actually forced}

Define, for an instance, $\alpha_{\mathrm{inst}}$ as the least
$\alpha$ such that some cost-good routing satisfies
$\flow^{\mathcal{P}}(a) \le x(a) + \alpha\dmax$ everywhere.

\begin{proposition}[machine-verified]
For the instance above, $\alpha_{\mathrm{inst}} = 16/15$.
\end{proposition}

In absolute terms, a cost-preserving rounding of this instance needs 16 units of
slack on one arc where the conjecture allows 15: the conjecture is refuted by
one unit. Consequently the instance says nothing about the question the
literature identifies as the real target, namely whether a cost-good routing
always exists within $x + O(\dmax)$; the only known upper bound remains 2, for
planar graphs \cite{TVZ24}, and no finite bound is known in general. We stress
this because the distinction is easily lost: refuting the constant 1 and
refuting $O(\dmax)$ are very different statements, and only the former has
happened.

\section{Consistency with everything proved, and the class boundary}

A counterexample must contradict no theorem. We tested each proved result
directly on the instance; all pass.

\begin{center}
\begin{tabular}{@{}lll@{}}
\toprule
Result & Test & Outcome\\
\midrule
Dinitz--Garg--Goemans \cite{DGG99} & a congestion-good routing exists & 4 exist\\
Morell--Skutella cost-free \cite{MS22} & some routing meets $x \pm \dmax$ & 4 witnesses\\
Conjecture 1.5 and planar \cite{TVZ24} & cost-good within $2\dmax$ & witnesses exist\\
Skutella \cite{Sku02} & demands not all multiples & $\{10, 15\}$\\
Majthoub Almoghrabi et al.\ \cite{MSW25} & not series-parallel & $K_4$ subdivision\\
Two-layer case \cite{TVZ24} & not a two-layer network & longest path 4 arcs\\
\bottomrule
\end{tabular}
\end{center}

Two consequences deserve statement. First, since the digraph is acyclic
(verified), any routing satisfying the Morell--Skutella cost conjecture
(Conjecture 1.4 of \cite{STVZ25}) would satisfy Goemans' conclusion, so that
conjecture is false as well; and if $x$ were a convex combination of
congestion-good unsplittable loads $y_i$, then $\min_i c^{\mathsf T} y_i \le
c^{\mathsf T} x$ would give a congestion-good routing of cost at most 58, so the
convex-combination form \cite{MSS07,MSW25} is false too. The cost-free
conjecture (Conjecture 1.3) is untouched, and with it Theorem 1.6 of
\cite{STVZ25}, which makes the cost-free statement the live route to a $2\dmax$
cost result.

Second, the branch vertices $\{s, u, v, w\}$ carry six internally disjoint
connecting paths, so the underlying graph contains a $K_4$ subdivision.
Series-parallel digraphs are $K_4$-minor-free, and the conjecture is proved for
them \cite{MSW25}. The counterexample therefore sits exactly one structure past
the class where the statement holds. Moreover the graph has no vertex of degree
4 or more and only four vertices of degree 3, so by Kuratowski's theorem it is
planar with no embedding computation; since \cite{TVZ24} proves the cost
statement for planar graphs at $2\dmax$, the planar constant lies strictly
between 1 and 2.

\section{The separation LP, and the optimality of the published prices}

Let $\mathcal{U}(x)$ be the set of loads of congestion-good routings: finite,
nonempty by \cite{DGG99}, and independent of $c$. An instance admits some
nonnegative cost vector making it a counterexample precisely when
\[
\text{(SEP)}\qquad
\max\ \delta \quad \text{s.t.} \quad
c^{\mathsf T}(y - x) \ge \delta \ \ \forall y \in \mathcal{U}(x), \qquad
\textstyle\sum_a c_a = 1, \qquad c \ge 0
\]
has optimum $\delta^\star > 0$. (If $\delta^\star > 0$ with witness $c$, every
congestion-good routing costs more than $x$; conversely a witnessing $c \neq 0$
has $\min_y c^{\mathsf T}(y-x) > 0$ over a finite nonempty set, and rescaling by
$1/\sum_a c_a$ gives a feasible point with positive objective.) We solve (SEP)
in exact rational arithmetic.

\begin{proposition}[machine-verified]
For the instance above, $\delta^\star = 2/7$, attained by
$c = \tfrac17(2, 3, 0, 2, 0, 0, 0, 0, 0)$, which is the published cost vector
normalised.
\end{proposition}

Hence the published prices are an optimal separator for that graph and
fractional flow: no reweighting of the arcs produces a larger normalised gap,
and the 58-versus-60 separation is the best that instance can do. Practically,
(SEP) removes the cost vector from any future search: an exhaustive hunt ranges
over graphs, demands and fractional flows only, with one exact linear program
per instance.

\section{Minimality, and a retraction}
\label{sec:minimality}

\begin{theorem}
No single-terminal instance is a counterexample, for any nonnegative cost
vector.
\end{theorem}

\begin{proof}
With $T = \{t\}$ and demand $d$ we have $\dmax = d$. Any routing loads its
chosen path with $d$ and the rest with 0, and $d \le x_a + d$ holds on every arc
because $x_a \ge 0$, so every routing is congestion-good. The flow $x$ has value
$d$ from $s$ to $t$, hence decomposes into $s$--$t$ paths of total weight $d$
plus circulations of nonnegative cost; each path costs at least the cheapest
path cost $\gamma$, so $c^{\mathsf T} x \ge d\gamma$, which is the cost of
routing on a cheapest path. That routing is therefore congestion-good and
cost-good.
\end{proof}

\begin{remark}[retraction]
In an earlier internal record we asserted, as a derivation, that no
counterexample exists with at most two terminals, on the grounds that a conflict
graph on at most two nodes has no odd cycle. That argument is invalid: a
terminal may have many path choices, so the conflict graph need not have two
nodes, and the step from an integral stable-set polytope to the absence of a
separating nonnegative cost vector was asserted rather than proved. We flagged
the claim as suspect in the declared hypothesis of the experiment that tested
it, and we withdraw it here. The two-terminal case is open.
\end{remark}

What does hold at two terminals, and in fact for any number of terminals in its
first half:

\begin{proposition}
Routing every terminal on a cheapest path is always cost-good. With two
terminals, that routing is congestion-good if and only if every arc lying on
both cheapest paths carries $x_a \ge \min(d_1, d_2)$. Consequently any
two-terminal counterexample contains an arc on both cheapest paths with
$x_a < \min(d_1, d_2)$.
\end{proposition}

\begin{proof}
Cost-goodness follows by applying the decomposition argument of the theorem to
each per-terminal flow in a decomposition $x = \sum_i x^i$. For the second part,
an arc used by exactly one of the two paths carries at most $\dmax \le x_a +
\dmax$; a shared arc carries $d_1 + d_2$, and $d_1 + d_2 \le x_a + \max(d_1,
d_2)$ is exactly $x_a \ge \min(d_1, d_2)$.
\end{proof}

We swept 184 two-terminal instances with a shared spine and two choices per
terminal, over four demand pairs, every integer split of each demand, and four
cost pairs, solving (SEP) exactly on each: none admits any nonnegative cost
vector making it a counterexample, and the characterisation above held on every
one. This is evidence over a bounded box, not a proof, and we state it as such.

\section{What remains open}

The refutation settles the constant 1 and leaves the quantitative question
essentially untouched. Writing $\alpha^\star$ for the least $\alpha$ such that
every instance admits a cost-good routing within $x + \alpha\dmax$, the record
after this work is $\alpha^\star \ge 16/15$ and, for planar instances,
$\alpha^\star_{\mathrm{planar}} \le 2$; in general no finite upper bound is
known. Three concrete questions follow. Can the conflict-triangle mechanism be
extended to longer odd cycles or to conflict cliques, whose stable-set gaps are
larger, and does flow structure permit the corresponding instances? What is the
minimum size of a counterexample in terminals, vertices and arcs? And does the
equality $\sum_i \rho_i = \alpha_{\mathrm{inst}} = 16/15$ observed here reflect
a general relation between the stable-set violation and the forced violation
budget, or is it a coincidence of this instance's calibration? We record it as
undecided.

\section*{Provenance and scope}

The counterexample instance is not ours. It was announced publicly on 22--23
July 2026, attributed to D.\ Rybin working with a large language model, and was
supplied to this programme as a file bundle which we archived with SHA-256
hashes. What is ours is the independent exact verification, the quantitative
constant $16/15$, the consistency battery against the proved results, the $K_4$
and planarity placements, the separation-LP optimality statement, the
single-terminal theorem, the two-terminal characterisation, and the retraction
recorded above. We make no claim about priority or attribution, which are not
mathematical facts, and we take no position on how the community will receive
the announcement. Every computation is exact and reproducible from the committed
code; hypotheses were declared and committed before each run, and the verdicts,
including the refuted prediction, are part of the public record.

\begin{thebibliography}{9}

\bibitem{DGG99} Y.\ Dinitz, N.\ Garg, M.\ X.\ Goemans,
\emph{On the single-source unsplittable flow problem},
Combinatorica \textbf{19}(1) (1999) 17--41.
\doi{10.1007/s004930050043}

\bibitem{Sku02} M.\ Skutella,
\emph{Approximating the single source unsplittable min-cost flow problem},
Mathematical Programming \textbf{91} (2002).

\bibitem{MSS07} M.\ Martens, F.\ Salazar, M.\ Skutella,
\emph{Convex combinations of single source unsplittable flows},
ESA 2007. \href{https://doi.org/10.1007/978-3-540-75520-3\_36}{doi:10.1007/978-3-540-75520-3\_36}

\bibitem{MS22} A.\ Morell, M.\ Skutella,
\emph{Single source unsplittable flows with arc-wise lower and upper bounds},
Mathematical Programming (2022).

\bibitem{TVZ24} V.\ Traub, L.\ Vargas Koch, R.\ Zenklusen,
\emph{Single-source unsplittable flows in planar and bounded-genus graphs},
Mathematical Programming (2026); preprint arXiv:2308.02651.
\doi{10.1007/s10107-026-02365-x}

\bibitem{MSW25} M.\ Majthoub Almoghrabi, M.\ Skutella, P.\ Warode,
\emph{Integer and unsplittable multiflows in series-parallel digraphs},
IPCO 2025; Mathematical Programming (2026); preprint arXiv:2412.05182.
\doi{10.1007/s10107-026-02392-8}

\bibitem{STVZ25} C.\ Swamy, V.\ Traub, L.\ Vargas Koch, R.\ Zenklusen,
\emph{Unsplittable cost flows from unweighted error-bounded variants},
preprint arXiv:2510.21287 (2025).
\doi{10.48550/arXiv.2510.21287}

\end{thebibliography}

\end{document}
